% Options for packages loaded elsewhere
\PassOptionsToPackage{unicode}{hyperref}
\PassOptionsToPackage{hyphens}{url}
\PassOptionsToPackage{dvipsnames,svgnames,x11names}{xcolor}
\documentclass[
  11pt,
]{article}
\usepackage{xcolor}
\usepackage[margin=1in]{geometry}
\usepackage{amsmath,amssymb}
\setcounter{secnumdepth}{-\maxdimen} % remove section numbering
\usepackage{iftex}
\ifPDFTeX
  \usepackage[T1]{fontenc}
  \usepackage[utf8]{inputenc}
  \usepackage{textcomp} % provide euro and other symbols
\else % if luatex or xetex
  \usepackage{unicode-math} % this also loads fontspec
  \defaultfontfeatures{Scale=MatchLowercase}
  \defaultfontfeatures[\rmfamily]{Ligatures=TeX,Scale=1}
\fi
\usepackage{lmodern}
\ifPDFTeX\else
  % xetex/luatex font selection
\fi
% Use upquote if available, for straight quotes in verbatim environments
\IfFileExists{upquote.sty}{\usepackage{upquote}}{}
\IfFileExists{microtype.sty}{% use microtype if available
  \usepackage[]{microtype}
  \UseMicrotypeSet[protrusion]{basicmath} % disable protrusion for tt fonts
}{}
\makeatletter
\@ifundefined{KOMAClassName}{% if non-KOMA class
  \IfFileExists{parskip.sty}{%
    \usepackage{parskip}
  }{% else
    \setlength{\parindent}{0pt}
    \setlength{\parskip}{6pt plus 2pt minus 1pt}}
}{% if KOMA class
  \KOMAoptions{parskip=half}}
\makeatother
\usepackage{longtable,booktabs,array}
\newcounter{none} % for unnumbered tables
\usepackage{calc} % for calculating minipage widths
% Correct order of tables after \paragraph or \subparagraph
\usepackage{etoolbox}
\makeatletter
\patchcmd\longtable{\par}{\if@noskipsec\mbox{}\fi\par}{}{}
\makeatother
% Allow footnotes in longtable head/foot
\IfFileExists{footnotehyper.sty}{\usepackage{footnotehyper}}{\usepackage{footnote}}
\makesavenoteenv{longtable}
\usepackage{graphicx}
\makeatletter
\newsavebox\pandoc@box
\newcommand*\pandocbounded[1]{% scales image to fit in text height/width
  \sbox\pandoc@box{#1}%
  \Gscale@div\@tempa{\textheight}{\dimexpr\ht\pandoc@box+\dp\pandoc@box\relax}%
  \Gscale@div\@tempb{\linewidth}{\wd\pandoc@box}%
  \ifdim\@tempb\p@<\@tempa\p@\let\@tempa\@tempb\fi% select the smaller of both
  \ifdim\@tempa\p@<\p@\scalebox{\@tempa}{\usebox\pandoc@box}%
  \else\usebox{\pandoc@box}%
  \fi%
}
% Set default figure placement to htbp
\def\fps@figure{htbp}
\makeatother
\setlength{\emergencystretch}{3em} % prevent overfull lines
\providecommand{\tightlist}{%
  \setlength{\itemsep}{0pt}\setlength{\parskip}{0pt}}
\usepackage{bookmark}
\IfFileExists{xurl.sty}{\usepackage{xurl}}{} % add URL line breaks if available
\urlstyle{same}
% fallback for those not using the hyperref driver hyperxmp:
\makeatletter
\@ifundefined{xmpquote}{\newcommand{\xmpquote}[1]{#1}}{}
\makeatother
\hypersetup{
  pdftitle={Aggregating RCTA Individual Forecasts},
  pdfauthor={Andy Hou},
  colorlinks=true,
  linkcolor={blue},
  filecolor={Maroon},
  citecolor={Blue},
  urlcolor={blue},
  pdfcreator={LaTeX via pandoc}}

\title{Aggregating RCTA Individual Forecasts}
\author{Andy Hou}
\date{}

\begin{document}
\maketitle

\section{Memo: Aggregating RCTA Individual
Forecasts}\label{memo-aggregating-rcta-individual-forecasts}

\subsection{Executive Summary}\label{executive-summary}

This memo analyzes individual forecasts from the first-year HFC RCTA
competition. The statistical problem is to construct, for each
question-day, a crowd probability vector from individual forecasts and
to compare aggregation rules by Brier loss. I reconstruct five baseline
aggregation methods from \texttt{rct-a-prediction-sets.csv}, score them
against outcomes from \texttt{rct-a-questions-answers.csv}, and then
propose a sixth method: \textbf{site-balanced bounded evidence pooling}.
The method clips individual probabilities to \texttt{{[}0.02,\ 0.98{]}},
pools clipped log-odds within forecasting site, and then combines sites
with equal weight. It improves over all five baselines under both
scoring summaries: question-balanced Brier decreases from 0.235399 for
the median baseline to 0.228379, and question-day weighted Brier
decreases from 0.218701 for geometric mean to 0.211749. The paired
bootstrap improvement over median is 0.0070198 with 95\% interval
\texttt{{[}-0.0003299,\ 0.0137224{]}}, so the effect is favorable but
modest. In practical terms, the method limits how much any one
forecaster can move the aggregate and treats each forecasting platform
as one information source rather than counting all of its members as
independent evidence.

\subsection{1. Data and Unit of
Analysis}\label{data-and-unit-of-analysis}

The raw RCTA data consist of three files. The questions-answers file
contains question metadata, answer options, dates, and resolved
outcomes. The prediction-sets file contains individual-level forecasts;
a prediction set is one forecaster's submission for one question, with
multiple rows for multi-option questions. The daily-forecasts file
contains performer-method forecasts rather than individual forecasts, so
it is not used as the source for aggregation.

The target unit is a question-day. Let \(q \in \{1,\ldots,Q\}\) index
questions, \(t \in \mathcal{T}_q\) index scoring days for question
\(q\), \(k \in \{1,\ldots,K_q\}\) index answer options, and
\(i \in \mathcal{I}_{qt}\) index forecasters active for question-day
\((q,t)\). Let \(s(i)\) denote the forecasting site of forecaster \(i\).
I write

\[
\mathcal{S}_{qt}=\{s(i): i\in\mathcal{I}_{qt}\}, \qquad
n_{sqt}=|\{i\in\mathcal{I}_{qt}:s(i)=s\}|
\]

for the set of sites and the number of forecasters from site \(s\) on
question-day \((q,t)\). The individual forecast vector is

\[
p_{iqt} = (p_{iqt1},\ldots,p_{iqtK_q}), \qquad \sum_{k=1}^{K_q}p_{iqtk}=1,
\]

and the resolved outcome vector is

\[
y_q = (y_{q1},\ldots,y_{qK_q}), \qquad y_{qk}\in\{0,1\}, \qquad \sum_{k=1}^{K_q}y_{qk}=1.
\]

For the 166 one-row binary questions, the stored answer option is
interpreted as the focal Yes answer. I construct

\[
p_{iqt,\mathrm{No}} = 1-p_{iqt,\mathrm{Yes}}, \qquad
y_{q,\mathrm{No}} = 1-y_{q,\mathrm{Yes}}.
\]

The cleaned scored sample contains 189 questions and 13,129
question-days.

\subsection{2. Cleaning and Scoring-Day
Construction}\label{cleaning-and-scoring-day-construction}

Forecast rows were retained only if they had non-missing and valid
forecast probabilities, non-missing resolved probabilities, valid
question and prediction-set identifiers, valid answer identifiers,
non-missing forecaster identifiers, non-missing timestamps, and were not
marked as made after correctness was known. Timestamps were converted to
America/New\_York. The scoring day is defined by the HFC 2:01pm ET
cutoff: forecasts before 2:01pm ET are assigned to that calendar day,
while forecasts at or after 2:01pm ET are assigned to the next scoring
day.

For each \((q,t,i)\), I selected the latest prediction set by timestamp
and retained all answer rows from that prediction set. This preserves
the vector-valued nature of multi-option submissions. Selecting answer
rows independently would create artificial forecast vectors that no
forecaster actually submitted.

\subsection{3. Loss Functions and
Estimands}\label{loss-functions-and-estimands}

For binary and nominal questions, the Brier loss for an aggregate
forecast \(\hat p_{qt}\) is

\[
L_{qt}(\hat p)=\sum_{k=1}^{K_q}(\hat p_{qtk}-y_{qk})^2.
\]

For binary questions this is the original 0-2 Brier scale:

\[
L_{qt}(\hat p)=
(\hat p_{\mathrm{Yes}}-y_{\mathrm{Yes}})^2+
(\hat p_{\mathrm{No}}-y_{\mathrm{No}})^2
=2(\hat p_{\mathrm{Yes}}-y_{\mathrm{Yes}})^2.
\]

For ordinal questions, I use the Jose-Nau-Winkler ordered Brier score.
Let

\[
F_j = \sum_{k\le j}\hat p_{qtk}, \qquad
O_j = \sum_{k\le j}y_{qk}, \qquad j=1,\ldots,K_q-1.
\]

The ordered score is

\[
L^{\mathrm{ord}}_{qt}(\hat p)=
\frac{2}{K_q-1}\sum_{j=1}^{K_q-1}(F_j-O_j)^2.
\]

I report two empirical risk summaries. The question-day weighted risk is

\[
\widehat R_{\mathrm{QD}}(m)=
\frac{1}{\sum_q |\mathcal{T}_q|}
\sum_{q=1}^{Q}\sum_{t\in\mathcal{T}_q}
L_{qt}(\hat p^{(m)}_{qt}),
\]

where each scored question-day receives one vote. The question-balanced
risk is

\[
\widehat R_{\mathrm{QB}}(m)=
\frac{1}{Q}\sum_{q=1}^{Q}
\left[
\frac{1}{|\mathcal{T}_q|}
\sum_{t\in\mathcal{T}_q}
L_{qt}(\hat p^{(m)}_{qt})
\right],
\]

where each question receives one vote regardless of how long it remained
open. The local HFC scoring document clearly establishes daily Brier /
MDB, but does not fully specify the second-stage aggregation for this
work-test comparison. I therefore report both estimands.

\subsection{4. Baseline Aggregation
Operators}\label{baseline-aggregation-operators}

For each question-day and answer option, I compute five baseline
aggregation rules. Let \(p_{ik}\) abbreviate \(p_{iqtk}\) within a fixed
\((q,t)\). The raw mean is

\[
a_k=\frac{1}{n}\sum_{i=1}^n p_{ik}, \qquad
\hat p_k=\frac{a_k}{\sum_j a_j}.
\]

The median is

\[
a_k=\mathrm{median}_{i}(p_{ik}), \qquad
\hat p_k=\frac{a_k}{\sum_j a_j}.
\]

The geometric mean of probabilities is

\[
a_k=\exp\left(\frac{1}{n}\sum_{i=1}^n\log(\mathrm{clip}(p_{ik},\epsilon,1-\epsilon))\right),
\qquad
\hat p_k=\frac{a_k}{\sum_j a_j}.
\]

The trimmed mean removes the top and bottom 10\% of scalar probabilities
before averaging. The geometric mean of odds is implemented as
one-vs-rest log-odds pooling:

\[
z_k=\frac{1}{n}\sum_{i=1}^n
\log\left[
\frac{\mathrm{clip}(p_{ik},\epsilon,1-\epsilon)}
{1-\mathrm{clip}(p_{ik},\epsilon,1-\epsilon)}
\right],
\qquad
a_k=\sigma(z_k),
\qquad
\hat p_k=\frac{a_k}{\sum_j a_j}.
\]

For binary questions, this reduces to

\[
\hat p_{\mathrm{Yes}}=\sigma\left(
\frac{1}{n}\sum_{i=1}^n \mathrm{logit}(p_{i,\mathrm{Yes}})
\right),
\qquad
\hat p_{\mathrm{No}}=1-\hat p_{\mathrm{Yes}}.
\]

The implementation uses \(\epsilon=10^{-6}\) for the original geometric
and odds-based baselines.

\subsection{5. Baseline Results}\label{baseline-results}

\begin{figure}
\centering
\pandocbounded{\includegraphics[keepaspectratio,alt={Baseline scores under two estimands}]{output/figure/baseline_two_metrics.png}}
\caption{Baseline scores under two estimands}
\end{figure}

{\def\LTcaptype{none} % do not increment counter
\begin{longtable}[]{@{}lrr@{}}
\toprule\noalign{}
Method & Question-Balanced Brier & Question-Day Weighted Brier \\
\midrule\noalign{}
\endhead
\bottomrule\noalign{}
\endlastfoot
Median & 0.235399 & 0.219796 \\
Trimmed mean & 0.237679 & 0.225603 \\
Raw mean & 0.243280 & 0.231857 \\
Geometric mean & 0.249192 & 0.218701 \\
Geometric mean of odds & 0.250367 & 0.233909 \\
\end{longtable}
}

The two estimands give different rankings. Under
\(\widehat R_{\mathrm{QD}}\), geometric mean is best. Under
\(\widehat R_{\mathrm{QB}}\), median is best. This reversal indicates a
robustness tradeoff. Geometric methods perform well on many scored days,
but they can incur large losses on a smaller number of questions. Median
is more stable at the question level, but it does not accumulate
evidence on a log scale.

\subsection{6. Diagnostic Analysis}\label{diagnostic-analysis}

Before proposing an improved aggregation rule, I diagnosed which failure
modes are supported by the data. The key diagnostic was the Murphy
decomposition, written on the standard 0-1 Brier scale as

\[
\mathrm{BS}=\mathrm{Reliability}-\mathrm{Resolution}+\mathrm{Uncertainty}.
\]

Reliability measures calibration error; resolution measures the ability
to separate events from non-events. For median, the reliability share is
0.008. For GMO, it is 0.026. Thus the leading baselines are not
dominated by calibration error.

{\def\LTcaptype{none} % do not increment counter
\begin{longtable}[]{@{}
  >{\raggedright\arraybackslash}p{(\linewidth - 4\tabcolsep) * \real{0.3333}}
  >{\raggedright\arraybackslash}p{(\linewidth - 4\tabcolsep) * \real{0.3333}}
  >{\raggedright\arraybackslash}p{(\linewidth - 4\tabcolsep) * \real{0.3333}}@{}}
\toprule\noalign{}
\begin{minipage}[b]{\linewidth}\raggedright
Diagnostic
\end{minipage} & \begin{minipage}[b]{\linewidth}\raggedright
Empirical result
\end{minipage} & \begin{minipage}[b]{\linewidth}\raggedright
Implication
\end{minipage} \\
\midrule\noalign{}
\endhead
\bottomrule\noalign{}
\endlastfoot
Murphy decomposition & median rel\_share = 0.008; GMO rel\_share = 0.026
& Global recalibration / extremization has little target \\
Type stratification & binary 0.147, nominal 0.487, ordinal 0.253 under
median question-day weighted scoring & Question geometry matters \\
Dispersion stratification & low dispersion 0.141, high dispersion 0.273
& Disagreement marks difficult question-days \\
Forecast age & same-day 0.219, one-day 0.220 & Much weaker than type and
dispersion in this no-carry-forward pipeline \\
Staleness interpretation & retained forecasts are same-day or one-day
latest submissions, not carried-forward stale forecasts & The pipeline
limits what recency decay can change \\
\end{longtable}
}

These diagnostics rule out the most standard improvement in the
forecasting literature: global extremization. Extremization is useful
when an average forecast is too close to 0.5 and should be pushed
outward. Here, the leading baselines are not mainly underconfident. The
problem is on the other side. Evidence-scale pooling is often useful,
but unbounded log-odds evidence can become too extreme.

This gives the final method a clear relation to the literature.
Extremization corrects under-accumulated evidence by moving an average
forecast away from 0.5. Bounded evidence pooling corrects
over-accumulated evidence by limiting how far one forecaster can move
the log-odds average. They are mirror-image operations on the same
calibration spectrum: one pushes a conservative aggregate outward, while
the other pulls an overconfident evidence pool back inward.

\subsection{7. Tested Conditional
Improvements}\label{tested-conditional-improvements}

I also tested the feature classes suggested by the prompt. The results
are summarized in Table 4.

{\def\LTcaptype{none} % do not increment counter
\begin{longtable}[]{@{}
  >{\raggedright\arraybackslash}p{(\linewidth - 6\tabcolsep) * \real{0.2500}}
  >{\raggedright\arraybackslash}p{(\linewidth - 6\tabcolsep) * \real{0.2500}}
  >{\raggedright\arraybackslash}p{(\linewidth - 6\tabcolsep) * \real{0.2500}}
  >{\raggedright\arraybackslash}p{(\linewidth - 6\tabcolsep) * \real{0.2500}}@{}}
\toprule\noalign{}
\begin{minipage}[b]{\linewidth}\raggedright
Candidate direction
\end{minipage} & \begin{minipage}[b]{\linewidth}\raggedright
Statistical object
\end{minipage} & \begin{minipage}[b]{\linewidth}\raggedright
Result
\end{minipage} & \begin{minipage}[b]{\linewidth}\raggedright
Decision
\end{minipage} \\
\midrule\noalign{}
\endhead
\bottomrule\noalign{}
\endlastfoot
Skill weighting & Prior resolved questions per active forecaster &
First-half median = 3.5; second-half median = 16 & Not stable enough for
season-wide use \\
Type-conditioned selection & Walk-forward best method by question type &
Oracle improves slightly; walk-forward underperforms median & Hindsight
structure is not deployable signal \\
Thin-crowd override & Switch method when \(n_{qt}\) is small & Best
threshold mostly collapses to median & Not a meaningful improvement \\
Dispersion shrinkage & Shrink high-dispersion forecasts toward uniform &
Brier worsens monotonically as shrinkage increases & Dispersion is a
marker, not an intervention \\
Activity/update history & Accuracy change after forecast updates & Large
updates improve 79\% of the time & Interesting future signal, not final
rule \\
\end{longtable}
}

The negative results are important. They show that the final method is
not a default fallback; it is chosen after plausible leakage-free
alternatives fail for identifiable statistical reasons.

\subsection{8. GMO Tail Autopsy}\label{gmo-tail-autopsy}

Let

\[
\Delta_{qt}=L_{qt}(\hat p^{\mathrm{GMO}})-L_{qt}(\hat p^{\mathrm{median}}).
\]

Negative \(\Delta_{qt}\) means GMO beats median on that question-day.
GMO beats median on 68.0\% of all question-days and 75.3\% of binary
question-days. However, positive excess loss is concentrated in a small
number of questions. In binary questions, the worst 10 questions account
for 85.2\% of GMO's positive excess loss relative to median, and the
worst 20 account for 97.0\%.

This is the empirical basis for the final method. The issue is not that
log-odds pooling is always poor. Rather, it is a high-variance
evidence-pooling rule: often useful, occasionally catastrophic. Since
the original GMO clips only at \(10^{-6}\), one near-certain forecast
can contribute approximately \(\pm 13.8\) log-odds units. That is an
implausibly large individual evidence contribution in geopolitical
forecasting.

The relative-to-individual diagnostic gives the positive side of the
same story. GMO's mean percentile rank among individual forecasters is
0.590 on binary question-days and 0.600 on ordinal question-days,
compared with 0.447 and 0.563 for median. So GMO often extracts crowd
signal better than the median. Its problem is not typical performance;
it is tail risk.

\subsection{9. Final Method: Site-Balanced Bounded Evidence
Pooling}\label{final-method-site-balanced-bounded-evidence-pooling}

The final method modifies GMO in two ways. First, it bounds individual
evidence. Second, it changes the aggregation unit from forecaster to
site.

In plain language, I first limit how much any one forecaster can
influence the aggregate by capping extreme probabilities. I then average
forecasts within each forecasting site and combine sites with equal
weight. This prevents one overconfident forecaster, or one large site,
from dominating the crowd forecast.

The method satisfies the prompt's temporal constraint directly. It uses
the forecasts available on day \(T\), time-invariant site labels, and a
fixed cap \(c=0.02\). The cap is not fit from resolved outcomes and does
not use information from day \(T\) or later.

For binary questions, set \(c=0.02\) and define

\[
p_i^{(c)}=\min\{\max(p_i,c),1-c\}.
\]

The individual bounded evidence is

\[
z_i^{(c)}=\mathrm{logit}(p_i^{(c)}).
\]

Within site \(s\), pool evidence by

\[
\bar z_s=\frac{1}{n_s}\sum_{i:s(i)=s} z_i^{(c)}.
\]

Across sites, combine site-level evidence by

\[
\bar z=\frac{1}{|\mathcal{S}_{qt}|}\sum_{s\in\mathcal{S}_{qt}}\bar z_s,
\qquad
\hat p_{\mathrm{Yes}}=\sigma(\bar z).
\]

For multi-option questions, I use the one-vs-rest extension. For each
option \(k\),

\[
\bar z_{sk}=\frac{1}{n_s}\sum_{i:s(i)=s}\mathrm{logit}(p_{ik}^{(c)}),
\qquad
\bar z_k=\frac{1}{|\mathcal{S}_{qt}|}\sum_s \bar z_{sk},
\qquad
a_k=\sigma(\bar z_k),
\qquad
\hat p_k=\frac{a_k}{\sum_j a_j}.
\]

The cap \(c=0.02\) implies that no individual forecast contributes more
than \(49:1\) odds. This is a modeling assumption, not merely a
numerical safeguard. It says that individual near-certainty is treated
as bounded evidence rather than an unbounded likelihood ratio.

The following sensitivity check varies the cap while keeping the same
site-balanced structure. These rows are produced directly by
\texttt{analysis.py}, using the same scoring code as the final method.

{\def\LTcaptype{none} % do not increment counter
\begin{longtable}[]{@{}rrr@{}}
\toprule\noalign{}
Cap \(c\) & Maximum Individual Odds & Question-Balanced Brier \\
\midrule\noalign{}
\endhead
\bottomrule\noalign{}
\endlastfoot
0.005 & 199:1 & 0.230249 \\
0.010 & 99:1 & 0.228757 \\
0.020 & 49:1 & 0.228379 \\
0.050 & 19:1 & 0.232418 \\
\end{longtable}
}

The site-level step is also substantive. HFC sites are hybrid
forecasting systems, not ordinary demographic groups. Forecasters within
a site may share platform design, information feeds, training,
discussion environment, and machine assistance. Treating every
forecaster as an independent evidence source can therefore over-count
large sites.

\subsection{10. Step 4 Results and
Ablation}\label{step-4-results-and-ablation}

\begin{figure}
\centering
\pandocbounded{\includegraphics[keepaspectratio,alt={Step 4 ablation}]{output/figure/step4_ablation.png}}
\caption{Step 4 ablation}
\end{figure}

{\def\LTcaptype{none} % do not increment counter
\begin{longtable}[]{@{}
  >{\raggedright\arraybackslash}p{(\linewidth - 6\tabcolsep) * \real{0.2143}}
  >{\raggedleft\arraybackslash}p{(\linewidth - 6\tabcolsep) * \real{0.2857}}
  >{\raggedleft\arraybackslash}p{(\linewidth - 6\tabcolsep) * \real{0.2857}}
  >{\raggedright\arraybackslash}p{(\linewidth - 6\tabcolsep) * \real{0.2143}}@{}}
\toprule\noalign{}
\begin{minipage}[b]{\linewidth}\raggedright
Method
\end{minipage} & \begin{minipage}[b]{\linewidth}\raggedleft
Question-Balanced Brier
\end{minipage} & \begin{minipage}[b]{\linewidth}\raggedleft
Question-Day Weighted Brier
\end{minipage} & \begin{minipage}[b]{\linewidth}\raggedright
Interpretation
\end{minipage} \\
\midrule\noalign{}
\endhead
\bottomrule\noalign{}
\endlastfoot
Median baseline & 0.235399 & 0.219796 & Robust baseline \\
Original GMO & 0.250367 & 0.233909 & Correct but fragile \\
Capped GMO 0.02 & 0.232327 & 0.216881 & Bounded evidence only \\
Site-balanced median & 0.234246 & 0.221594 & Site unit correction
only \\
Site-balanced capped GMO 0.02 & 0.228379 & 0.211749 & Combined method \\
\end{longtable}
}

The ablation supports the proposed mechanism. Bounding individual
evidence improves over median. Site balancing alone gives a smaller
improvement under question-balanced scoring. The combined method
performs best, suggesting that individual-level overconfidence and
site-level non-exchangeability are distinct failure modes.

By question type, the final method improves the median baseline under
question-balanced scoring:

{\def\LTcaptype{none} % do not increment counter
\begin{longtable}[]{@{}lrr@{}}
\toprule\noalign{}
Question Type & Median Baseline & Step 4 Method \\
\midrule\noalign{}
\endhead
\bottomrule\noalign{}
\endlastfoot
Binary & 0.189982 & 0.179451 \\
Nominal & 0.478601 & 0.457898 \\
Ordinal & 0.244888 & 0.243401 \\
\end{longtable}
}

The paired bootstrap over questions gives

\[
\widehat R_{\mathrm{QB}}(\mathrm{median})-\widehat R_{\mathrm{QB}}(\mathrm{Step4})
=0.0070198,
\]

with 95\% interval \texttt{{[}-0.0003299,\ 0.0137224{]}}. The point
estimate is favorable, but the interval slightly crosses zero. I
therefore interpret the improvement as promising and mechanistically
coherent, not as conclusive out-of-sample evidence.

\subsection{11. Temporal Validation}\label{temporal-validation}

\begin{figure}
\centering
\pandocbounded{\includegraphics[keepaspectratio,alt={Rolling-window validation}]{output/figure/rolling_window_validation.png}}
\caption{Rolling-window validation}
\end{figure}

Rolling-window checks ask whether the improvement is concentrated in a
short period. They are not external validation, but they are useful
temporal diagnostics.

{\def\LTcaptype{none} % do not increment counter
\begin{longtable}[]{@{}lrr@{}}
\toprule\noalign{}
Window & Mean Improvement & Share Positive Windows \\
\midrule\noalign{}
\endhead
\bottomrule\noalign{}
\endlastfoot
7-day daily-weighted & 0.005949 & 0.690608 \\
14-day daily-weighted & 0.006045 & 0.741379 \\
30-day daily-weighted & 0.006891 & 0.797468 \\
45-day daily-weighted & 0.007927 & 0.839161 \\
14-day question-balanced & 0.006045 & 0.741379 \\
30-day question-balanced & 0.006891 & 0.797468 \\
45-day question-balanced & 0.007927 & 0.839161 \\
\end{longtable}
}

The fixed method has positive mean improvement across all reported
windows and is positive in roughly 69-84\% of windows. This reduces the
concern that the improvement is driven by a single short episode.

\subsection{12. Limitations}\label{limitations}

The sample has only 189 scored questions, and the paired bootstrap
interval crosses zero. The result should therefore be validated on
another HFC season before being treated as a stable empirical law.

The cap \(c=0.02\) is interpretable, but still a modeling choice. The
sensitivity table shows that nearby caps give similar performance, but a
stronger design would pre-specify the cap or select it strictly by
walk-forward validation.

The site-balanced component has two possible mechanisms: correcting
within-site dependence and down-weighting weaker large sites. The data
support both as plausible. This is a limitation for causal
interpretation, although not necessarily for predictive performance.

The multi-option GMO extension is not unique. I use one-vs-rest odds
pooling followed by normalization. Other log-evidence geometries, such
as centered log-ratio pooling, may perform differently for nominal and
ordinal questions.

Finally, the analysis does not carry forecasts forward. This is
consistent with the question-day construction used here, but it limits
conclusions about long-horizon staleness and recency decay.

\subsection{Conclusion}\label{conclusion}

The main statistical finding is that the five baseline methods expose a
tradeoff between evidence pooling and robustness. Geometric methods are
strong under question-day weighting but vulnerable to per-question tail
failures. Median is robust but does not pool evidence on the log-odds
scale. Site-balanced bounded evidence pooling resolves this tradeoff by
retaining log-odds evidence aggregation while bounding individual
influence and correcting the unit of aggregation from forecaster to
forecasting site. The resulting method improves both reported risk
summaries and has a coherent statistical mechanism, though the size of
the improvement should be interpreted conservatively.

\end{document}
